\section{Scope}
\label{sec:scope}

This section is the baseline the later sections detail. Every item has an identifier that the API,
front-end, chart, test and schedule sections use: M for Must, S for Should, C for Could, W for
Won't. Views are V1 to V8, charts G1 to G9, endpoints E1 to E22, journeys J1 to J8 and tasks T01
to T52.

\subsection{Scope statement}

\textbf{The product.} A web application for one Bitcoin address at a time. It shows the address's
confirmed activity per day, lists each day's transactions, marks the transactions that are unusual
compared with the same address's own earlier activity, says why in plain words, shows what each of
nine models makes of a transaction, and lets the user test a hypothetical transaction against the
address's real history.

\textbf{The user.} A student researcher or analyst deciding which transactions deserve a closer
look, as in Assignment 1. The application is used by one person on their own machine.

\textbf{The success boundary}, unchanged from Assignment 1: the system can explain why observed
transactions stand out; it cannot establish wrongdoing, identity, intent or ownership.

\textbf{The topic in the brief's list.} Area 5, cybersecurity threat detection in banking and
finance, by way of anomaly detection on transaction histories (area 4's method). The tutor approved
the topic for all three assignments.

\subsection{User journeys}
\label{sec:journeys}

Each journey is a path the video shows and an acceptance test runs (Section~\ref{sec:acceptance}).

\begin{longtable}{@{}L{0.8cm}L{4.2cm}L{6.9cm}L{2.1cm}@{}}
\toprule
\textbf{ID} & \textbf{Journey} & \textbf{Steps} & \textbf{Serves} \\
\midrule
\endhead
J1 & Open an address and read its activity & Pick an address from the list → the overview shows the date range, the source and snapshot date, four summary tiles, the daily chart and the days with unusual transactions & FR-01, FR-02, AC-01 \\
J2 & Find out why a transaction stands out & Click a flagged day → the day's transactions → one transaction: its facts, the reason in plain words, the comparison basis, the limitation line, the feature chart and every model's verdict & FR-03, FR-04, AC-03 \\
J3 & Type an address that cannot be shown & A typo is caught before any request; a valid address outside the snapshot gets an explanation and three supported examples; a short-history address lists its transactions unscored & FR-05, AC-02 \\
J4 & Test a hypothetical transaction & Open the what-if lab → enter an amount, a direction and a time → the verdict and the response curve update as the values change → keep the scenario & Brief: user input → prediction; M-03, M-04, M-06 \\
J5 & Explore with the charts & Zoom the daily chart with the range slider; filter the scatter to outgoing, flagged transactions; hover for the facts; click a point to open it & M-07, M-08 \\
J6 & Take the results away & Export the address's scored transactions as CSV; save a chart as PNG & M-12 \\
J7 & Compare the models and keep a watchlist & Switch the deciding method and watch the flags change; read the models page; add the address to the watchlist with a note; find it again from the first screen & S-01 to S-04 \\
J8 & Look up an address outside the snapshot & Start a live lookup → watch the progress → the same overview opens for the fetched address, marked as live data & S-09 (CR-08) \\
\bottomrule
\end{longtable}

\subsection{Scope table}
\label{sec:scopetable}

\must{} items together reach the full-points description of every rubric line.
\should{} items are the High Distinction extras and the polish; each stands alone and can be cut
without breaking another. \could{} items are built only after MS3 and only with a logged change.
\wont{} items are excluded, with the reason.

{\small
\begin{longtable}{@{}L{0.9cm}L{6.9cm}L{3.6cm}L{2.5cm}@{}}
\toprule
\textbf{ID} & \textbf{Item} & \textbf{Rubric lines it serves} & \textbf{Specified in} \\
\midrule
\endhead
\multicolumn{4}{@{}l}{\must} \\
M-01 & Back-end service: application factory, models loaded once at start, settings, CORS, logging, request identifiers, one error body, health check, generated API documentation & Code 1, 3; Functionality 4 & \ref{sec:architecture}, \ref{sec:api} \\
M-02 & Assignment 1's five endpoints on the real snapshot (E3 to E7), with its 400, 404 and 422 bodies & Functionality 1, 2, 4 & \ref{sec:api} \\
M-03 & What-if prediction (E8): a hypothetical transaction scored by every model that can score new input & Functionality 1, 2 & \ref{sec:whatif} \\
M-04 & What-if response curve (E9): one request scores a range of amounts & Functionality 1, 5; Visualisation 1 & \ref{sec:whatif} \\
M-05 & The three prototype screens on the real API (V1 to V3) through an Axios client; prototype banner and tools removed & Functionality 2; Usability 1 & \ref{sec:frontend} \\
M-06 & The what-if lab (V4): a validated form with live feedback & Functionality 2, 3; Usability 1, 2 & \ref{sec:viewwhatif} \\
M-07 & Charts G1, G2, G3 and G5 in Plotly.js under one chart contract: zoom, pan, range slider, tooltips, legend, units, snapshot date, table view, click-through & Visualisation 1 to 4 & \ref{sec:charts} \\
M-08 & A filter bar (status, direction, date range) that drives the scatter and the transaction table through E6 & Visualisation 2; Functionality 2 & \ref{sec:filters} \\
M-09 & Validation on both tiers from one rule table (address, date, transaction ID, what-if, list parameters), with shared test vectors & Functionality 3 & \ref{sec:validation} \\
M-10 & Error handling on both tiers: one catalogue of states, each with a next action; error boundary; toasts; retry; time limits; service-status indicator & Functionality 3; Usability 2 & \ref{sec:errors} \\
M-11 & Real-time behaviour: the what-if refreshes as values change; feedback within 100 ms on every request; layouts at 390, 820 and 1280 pixels & Functionality 5; Usability 1 & \ref{sec:realtime} \\
M-12 & Export: an address's scored transactions as CSV or JSON (E10); any chart as PNG & Usability 2; HD list & \ref{sec:api}, \ref{sec:chartcontract} \\
M-13 & Every model's verdict on a transaction (a table in V3), the method list (E11) and the models and limits page (V5) & Functionality 1 & \ref{sec:model}, \ref{sec:frontend} \\
M-14 & Tests on both tiers, the wording test, the README, and a clean install by a teammate on Windows and macOS & Code 1 to 3 & \ref{sec:quality} \\
M-15 & The usability walkthrough (four tasks, two or more participants) and the accessibility inspection, with every severity 3 or 4 issue fixed & Usability 1, 2 & \ref{sec:usability} \\
M-16 & The submission: report, video, minutes, contribution form, zip & Report 1 to 3; all & \ref{sec:report} \\
\midrule
\multicolumn{4}{@{}l}{\should{} (in build order)} \\
S-01 & Method selector: view the flags of any of the nine methods (\code{?method=}); the baseline is the default & Functionality 1; Visualisation 2 & \ref{sec:methods} \\
S-02 & Chart G4: each model's score against its own threshold for one transaction & Visualisation 1, 2 & \ref{sec:charts} \\
S-03 & Chart G6: each method's recall on planted test transactions, from Assignment 2's tables, on V5 & Visualisation 1, 3 & \ref{sec:charts} \\
S-04 & Watchlist with a label and a note (E13 to E15), recent addresses, remembered preferences & Functionality 4; Usability 2 & \ref{sec:viewwatch} \\
S-05 & Explorer Assistant explanations on request (E12): guarded, cached, optional weights, prepared answers for the snapshot's flags & Functionality 1 & \ref{sec:assistant} \\
S-06 & Replay: an address's history played forward, the charts updating as each transaction arrives & Functionality 5; Visualisation 2 & \ref{sec:realtime} \\
S-07 & Linked selection: points selected in the scatter filter the table and show a summary & Visualisation 2 & \ref{sec:charts} \\
S-08 & Scenario table in the what-if lab: keep, compare and export what-ifs & Usability 2 & \ref{sec:viewwhatif} \\
S-09 & Live lookup of any valid Bitcoin address (E16 to E20): progress stream, cache, limits. Needs CR-08 & Functionality 1, 5 & \ref{sec:live} \\
S-10 & TypeScript types generated from the OpenAPI schema; tests and builds run on every push & Code 3 & \ref{sec:tests} \\
S-11 & Browser tests of the journeys at three widths, producing the report's screenshots & Code 1; Usability 1 & \ref{sec:tests} \\
S-12 & A script that measures response times, bundle sizes and chart render times, for the report & Visualisation 4 & \ref{sec:budgets} \\
\midrule
\multicolumn{4}{@{}l}{\could} \\
C-01 & Chart G7: TxGPT's surprisal per field for a snapshot transaction (E21) & Visualisation 1 & \ref{sec:charts} \\
C-02 & Activity profile card from Assignment 2's address profiles (E22) & Functionality 1 & \ref{sec:api} \\
C-03 & Charts G8 (calendar heat map) and G9 (snapshot overview) & Visualisation 1 & \ref{sec:charts} \\
C-04 & Review marks on transactions (``worth a look'', ``as expected'', ``unsure'') with a note and an export & Usability 2 & -- \\
C-05 & One-command mode: FastAPI serves the built front end; a Docker Compose file & Code 1 & \ref{sec:runmodes} \\
C-06 & Watch mode for a live address: new confirmed transactions appear by themselves & Functionality 5 & \ref{sec:live} \\
C-07 & A printable one-page summary of an address & Usability 2 & -- \\
\midrule
\multicolumn{4}{@{}l}{\wont} \\
W-01 & Any claim about who owns an address, links between addresses, or wallet tracking & \multicolumn{2}{L{6.1cm}}{Assignment 1 exclusion; the data cannot support it} \\
W-02 & Other blockchains and Bitcoin test networks & \multicolumn{2}{L{6.1cm}}{Assignment 1 exclusion; validation rejects them with a reason} \\
W-03 & Accounts, sign-in, roles & \multicolumn{2}{L{6.1cm}}{Not in the brief; one local user (D-A3-8)} \\
W-04 & The Elliptic and BitcoinHeist classifiers in the interface & \multicolumn{2}{L{6.1cm}}{Elliptic's licence forbids shipping a derived model; the BitcoinHeist model is 276 MB and near chance across time; a form for them would read as a verdict on an address} \\
W-05 & Free chat with the 152M model & \multicolumn{2}{L{6.1cm}}{Its open answers are not guarded; the glossary it learned has one outdated sentence (Section~\ref{sec:looseends})} \\
W-06 & Retraining or re-tuning any model, or changing an Assignment 2 number & \multicolumn{2}{L{6.1cm}}{The models are frozen at \code{results-v1.1}} \\
W-07 & Scoring an uploaded file of transactions & \multicolumn{2}{L{6.1cm}}{The live lookup covers new addresses; the address store (Section~\ref{sec:store}) leaves room for it later} \\
W-08 & A hosted public deployment as a marked item & \multicolumn{2}{L{6.1cm}}{The brief marks a local run from the README; asked of the tutor (Section~\ref{sec:tutorq})} \\
\bottomrule
\end{longtable}
}

\subsection{The minimum viable demonstration}
\label{sec:mvd}

Assignment 1 defined one complete journey as the gate before anything optional. Assignment 3's gate
is that journey on real data plus the brief's own minimum, and it is milestone MS2.

\begin{done}
On all 28 snapshot addresses, with both servers started from the README:
\begin{enumerate}
  \item J1, J2 and J3 pass: AC-01, AC-02, AC-03 and all five FR-05 states, checked by the
    acceptance tests and by hand on three addresses.
  \item J4 passes: a what-if entered in the form returns every model's verdict, and an invalid
    value is refused in the browser and again by the server.
  \item Three chart types are on screen (G1 bar, G2 scatter, G5 line) with zoom and tooltips,
    each with a title, axis units, a legend and the snapshot date.
  \item The API uses GET and POST, and its documentation page lists every endpoint.
  \item The wording test finds none of the five banned words in any response or on any screen,
    and every flagged transaction shows the limitation line.
  \item The layouts pass at 390, 820 and 1280 pixels with no sideways scrolling.
\end{enumerate}
\end{done}

\subsection{Cut order if time runs short}
\label{sec:cutorder}

Items are cut from the bottom of this list upwards, and each cut is recorded in the minutes and the
change log. A Should item that is not finished and tested by MS3 (Fri 27 Nov) is cut, not carried.

\begin{enumerate}
  \item Every Could item (never started before MS3).
  \item S-12, S-11, S-10: evidence and tooling. The Must tests still run by hand.
  \item S-09 live lookup, the largest Should item and the only one that needs the network.
  \item S-08 scenario table, S-07 linked selection.
  \item S-06 replay.
  \item S-05 Assistant: the template explanation already satisfies FR-04.
  \item S-04 watchlist: the API then has GET and POST only, which still meets the brief.
  \item S-03 and S-02: the charts fall back to the four Must charts, which still cover three types.
  \item S-01 method selector: the verdict table of M-13 still shows all nine models.
\end{enumerate}

Nothing in the Must tier is cut. If a Must item is at risk, the team re-plans at once under
Assignment 1's variance rules (Section~\ref{sec:variance}).
